\section{Calculation of some multiplicities}
\label{section-calculation}

\noindent
To prove equality of certain cycles later on we need to
compute some multiplicities. Our main tool, besides the
elementary lemmas on multiplicities given in the previous section,
will be Algebra, Lemma \ref{algebra-lemma-order-vanishing-determinant}.

\begin{lemma}
\label{lemma-length-multiplication}
Let $R$ be a Noetherian local ring.
Let $M$ be a finite $R$-module. Let $x \in R$. Assume that
\begin{enumerate}
\item $\dim(\text{Supp}(M)) \leq 1$, and
\item $\dim(\text{Supp}(M/xM)) \leq 0$.
\end{enumerate}
Write
$\text{Supp}(M) = \{\mathfrak m, \mathfrak q_1, \ldots, \mathfrak q_t\}$.
Then
$$
e_R(M, 0, x) =
\sum\nolimits_{i = 1, \ldots, t}
\text{ord}_{R/\mathfrak q_i}(x)
\text{length}_{R_{\mathfrak q_i}}(M_{\mathfrak q_i}).
$$
\end{lemma}

\begin{proof}
We first make some preparatory remarks.
The result of the lemma holds if $M$ has finite length, i.e., if $t = 0$,
because both the left hand side and the right hand side are zero
in this case, see Lemma \ref{lemma-finite-periodic-length}.
Also, if we have a short exact sequence $0 \to M \to M' \to M'' \to 0$
of modules satisfying (1) and (2), then lemma for 2 out of 3
of these implies the lemma for the third by the
additivity of length (Algebra, Lemma \ref{algebra-lemma-length-additive}) and
additivty of multiplicities (Lemma \ref{lemma-additivity-periodic-length}).

\medskip\noindent
Denote $M_i$ the image of $M$ in $M_{\mathfrak q_i}$, so
$\text{Supp}(M_i) = \{\mathfrak m, \mathfrak q_i\}$.
The kernel and cokernel of the map $M \to \bigoplus M_i$
have support $\{\mathfrak m\}$ and hence have finite length.
By our preparatory remarks, it follows that it suffices to
prove the lemma for each $M_i$. Thus we may assume that
$\text{Supp}(M) = \{\mathfrak m, \mathfrak q\}$.
In this case we have a finite filtration
$M \supset \mathfrak qM \supset \mathfrak q^2M \supset \ldots \supset
\mathfrak q^nM = 0$ by Algebra, Lemma
\ref{algebra-lemma-Noetherian-power-ideal-kills-module}.
Again additivity shows that it suffices to prove the lemma
in the case $M$ is annihilated by $\mathfrak q$.
In this case we can view $M$ as a $R/\mathfrak q$-module,
i.e., we may assume that $R$ is a Noetherian local domain
of dimension $1$ with fraction field $K$.
Dividing by the torsion submodule, i.e., by the
kernel of $M \to M \otimes_R K = V$ (the torsion has
finite length hence is handled by our preliminary remarks)
we may assume that $M \subset V$ is a lattice
(Algebra, Definition \ref{algebra-definition-lattice}).
Then $x : M \to M$ is injective and
$\text{length}_R(M/xM) = d(M, xM)$
(Algebra, Definition \ref{algebra-definition-distance}). Since
$\text{length}_K(V) = \dim_K(V)$
we see that $\det(x : V \to V) = x^{\dim_K(V)}$ and
$\text{ord}_R(\det(x : V \to V)) = \dim_K(V) \text{ord}_R(x)$.
Thus the desired equality follows from
Algebra, Lemma \ref{algebra-lemma-order-vanishing-determinant}
in this case.
\end{proof}

\begin{lemma}
\label{lemma-additivity-divisors-restricted}
Let $R$ be a Noetherian local ring.
Let $x \in R$. If $M$ is a finite Cohen-Macaulay module over $R$
with $\dim(\text{Supp}(M)) = 1$ and $\dim(\text{Supp}(M/xM)) = 0$, then
$$
\text{length}_R(M/xM)
=
\sum\nolimits_i \text{length}_R(R/(x, \mathfrak q_i))
\text{length}_{R_{\mathfrak q_i}}(M_{\mathfrak q_i}).
$$
where $\mathfrak q_1, \ldots, \mathfrak q_t$ are the
minimal primes of the support of $M$. If $I \subset R$ is an ideal
such that $x$ is a nonzerodivisor on $R/I$ and $\dim(R/I) = 1$, then
$$
\text{length}_R(R/(x, I))
=
\sum\nolimits_i \text{length}_R(R/(x, \mathfrak q_i))
\text{length}_{R_{\mathfrak q_i}}((R/I)_{\mathfrak q_i})
$$
where $\mathfrak q_1, \ldots, \mathfrak q_n$ are the minimal
primes over $I$.
\end{lemma}

\begin{proof}
These are special cases of Lemma \ref{lemma-length-multiplication}.
\end{proof}

\noindent
Here is another case where we can determine the value of a multiplicity.

\begin{lemma}
\label{lemma-powers-period-length-zero}
Let $R$ be a ring. Let $M$ be an $R$-module.
Let $\varphi : M \to M$ be an endomorphism and $n > 0$
such that $\varphi^n = 0$ and such that $\Ker(\varphi)/\Im(\varphi^{n - 1})$
has finite length as an $R$-module.
Then
$$
e_R(M, \varphi^i, \varphi^{n - i}) = 0
$$
for $i = 0, \ldots, n$.
\end{lemma}

\begin{proof}
The cases $i = 0, n$ are trivial as $\varphi^0 = \text{id}_M$ by convention.
Let us think of $M$ as an $R[t]$-module where multiplication by $t$
is given by $\varphi$. Let us write
$K_i = \Ker(t^i : M \to M)$ and
$$
a_i = \text{length}_R(K_i/t^{n - i}M),\quad
b_i = \text{length}_R(K_i/tK_{i + 1}),\quad
c_i = \text{length}_R(K_1/t^iK_{i + 1})
$$
Boundary values are $a_0 = a_n = b_0 = c_0 = 0$.
The $c_i$ are integers for $i < n$ as $K_1/t^iK_{i + 1}$
is a quotient of $K_1/t^{n - 1}M$ which is assumed to have finite length.
We will use frequently that $K_i \cap t^jM = t^jK_{i + j}$.
For $0 < i < n - 1$ we have an exact sequence
$$
0 \to
K_1/t^{n - i - 1}K_{n - i} \to
K_{i + 1}/t^{n - i - 1}M \xrightarrow{t} K_i/t^{n - i}M
\to K_i/tK_{i + 1} \to 0
$$
By induction on $i$ we conclude that $a_i$ and $b_i$ are
integers for $i < n$ and that
$$
c_{n - i - 1} - a_{i + 1} + a_i - b_i = 0
$$
For $0 < i < n - 1$ there is a short exact sequence
$$
0 \to
K_i/tK_{i + 1} \to
K_{i + 1}/tK_{i + 2} \xrightarrow{t^i}
K_1/t^{i + 1}K_{i + 2} \to
K_1/t^iK_{i + 1} \to 0
$$
which gives
$$
b_i - b_{i + 1} + c_{i + 1} - c_i = 0
$$
Since $b_0 = c_0$ we conclude that $b_i = c_i$ for $i < n$.
Then we see that
$$
a_2 = a_1 + b_{n - 2} - b_1,\quad
a_3 = a_2 + b_{n - 3} - b_2,\quad \ldots
$$
It is straightforward to see that this implies $a_i = a_{n - i}$ as desired.
\end{proof}

\begin{lemma}
\label{lemma-multiply-period-length}
Let $(R, \mathfrak m)$ be a Noetherian local ring. Let
$(M, \varphi, \psi)$ be a $(2, 1)$-periodic complex over $R$
with $M$ finite and with cohomology groups of finite length over $R$.
Let $x \in R$ be such that $\dim(\text{Supp}(M/xM)) \leq 0$. Then
$$
e_R(M, x\varphi, \psi) = e_R(M, \varphi, \psi) - e_R(\Im(\varphi), 0, x)
$$
and
$$
e_R(M, \varphi, x\psi) = e_R(M, \varphi, \psi) + e_R(\Im(\psi), 0, x)
$$
\end{lemma}

\begin{proof}
We will only prove the first formula as the second is proved
in exactly the same manner.
Let $M' = M[x^\infty]$ be the $x$-power torsion submodule of $M$.
Consider the short exact sequence $0 \to M' \to M \to M'' \to 0$.
Then $M''$ is $x$-power torsion free (More on Algebra, Lemma
\ref{more-algebra-lemma-divide-by-torsion}).
Since $\varphi$, $\psi$ map $M'$ into $M'$
we obtain a short exact sequence
$$
0 \to (M', \varphi', \psi') \to (M, \varphi, \psi) \to
(M'', \varphi'', \psi'') \to 0
$$
of $(2, 1)$-periodic complexes. Also, we get a short exact sequence
$0 \to M' \cap \Im(\varphi) \to \Im(\varphi) \to \Im(\varphi'') \to 0$.
We have
$e_R(M', \varphi, \psi) = e_R(M', x\varphi, \psi) =
e_R(M' \cap \Im(\varphi), 0, x) = 0$
by Lemma \ref{lemma-compare-periodic-lengths}.
By additivity (Lemma \ref{lemma-additivity-periodic-length})
we see that it suffices to prove the lemma for $(M'', \varphi'', \psi'')$.
This reduces us to the case discussed in the next paragraph.

\medskip\noindent
Assume $x : M \to M$ is injective.
In this case $\Ker(x\varphi) = \Ker(\varphi)$.
On the other hand we have a short exact sequence
$$
0 \to \Im(\varphi)/x\Im(\varphi) \to
\Ker(\psi)/\Im(x\varphi) \to \Ker(\psi)/\Im(\varphi) \to 0
$$
This together with (\ref{equation-multiplicity-coker-ker}) proves the formula.
\end{proof}







